\section{Experiments and results}
\label{sec:experiments}

\subsection{Shared setup}
\label{sec:setup}

Two settings carry every result below, and both are described here once rather than
repeated per experiment. The \emph{synthetic} setting is a length-8 sequence over a 4-symbol vocabulary, so all
%% results/WS2_synthetic_ground_truth.json (task.n_sequences, task.regime_S_feasible_count)
$65{,}536$ sequences enumerate and every distribution is computed rather than sampled: no
divergence reported on this task is an estimate. A per-symbol multiplicity budget defines
the support constraint, satisfied by $36{,}120$ sequences, and each symbol carries a
charge; the expectation constraint asks for a mean total charge of zero. The \emph{crystal}
setting is the MP-20 subset of the Materials Project, serialized as text under the grammar
of \cref{sec:method}, where the constrained quantity is mass density in
\unit{\gram\per\centi\metre\cubed}.

Four training arms recur throughout, differing only in what they train through rather than
in architecture or mask. The \emph{mask-only} arm trains with the grammar mask and no
projection. The \emph{projection} arm trains through \cref{eq:projection} and is read
either without the layer, giving its backbone, or with it, giving the deployed system. The
\emph{sequence-level comparator} trains by maximum likelihood on exact samples from the
sequence-level tilt, which is distributional policy gradient's fixed point with unit
importance weights; that rule is chosen deliberately, because it gives the comparator the
best case of the method it stands for, so no gradient-estimator artefact can be blamed for
the outcome. Baselines needing no training are named where they are used.

Metrics follow the constraint. On the synthetic task we report the mean of the constrained
quantity, its variance, and the per-sequence violation rate, meaning the probability mass
on sequences whose own total misses the target. On the crystal task we report absolute
density error, structural validity under an independent parser, and unique-formula fraction
as a diversity measure. Every figure reported below is read from a stored artifact at build
time rather than transcribed.

\subsection{The surrogate targets a concentrated distribution}
\label{sec:concentration}

\paragraph{The comparator reaches its target}
%% results/WS7_a6_tilt_comparator.json (summary.kl_from_tilt, summary.mean_charge)
The comparator lands $0.0167 \pm 0.0048$ nats from the exact tilt with a mean charge of
$-0.062 \pm 0.122$ against a target of zero and no mass outside the mask. It is not a
failed run, which is what makes the comparison in \cref{tab:concentration} informative
rather than a win over something broken.

\begin{table}[t]
  \centering
  %% Every value: results/A2_training_summary.json (arms.*, [mean, sd] PAIRS -- index [0])
  %% and results/WS7_a6_tilt_comparator.json (summary.*, reference.tilt,
  %% reference.base_S_conditioned). Exact enumeration, no sampling.
  \caption{\textbf{Both arms satisfy the expectation constraint; only one concentrates
    it.} The sequence-level optimum's charge variance exceeds the unconstrained data's,
    while the projection's is far below it. Per-sequence violation is the mass on
    sequences whose own total charge is nonzero. Mean $\pm$ s.d.\ over seeds where an
    arm is trained; the tilt and the data distribution are exact.}
  \label{tab:concentration}
  \small
  \begin{tabular*}{\textwidth}{@{\extracolsep{\fill}}l S[table-format=+1.4] S[table-format=1.4] S[table-format=2.3]@{}}
    \toprule
    \kilth{Arm} & {\kilth{Mean charge}} & {\kilth{Violation}} & {\kilth{Variance}} \\
    \midrule
    Data, mask-conditioned (exact)   & +2.5164 & 0.9109 &  8.564 \\
    Mask only                        & +2.4694 & 0.9083 &  8.762 \\
    Sequence-level comparator        & -0.0619 & 0.8751 & 10.362 \\
    Tilt, sequence-level optimum     & +0.0000 & 0.8755 & 10.452 \\
    \midrule
    Projection (deployed)            & \best{+0.0170} & \best{0.6155} & \best{1.547} \\
    \bottomrule
  \end{tabular*}
\end{table}

\begin{figure}[t]
  \centering
  %% figures/fig1_charge_distributions.pdf <- results/WS7_charge_marginals.json
  %% Exact enumerated distributions marginalized onto integer total charge.
  \kilgraphics{figures/fig1_charge_distributions.pdf}
  \caption{\textbf{The optimum spreads the constrained quantity; the projection
    concentrates it.} Exact charge distributions over all sequences. The
    sequence-level optimum and the model trained toward it place mass at both extremes
    to buy a mean of zero, ending wider than the mask-only arm they started from. The
    projection instead spikes at the target. Legend entries give each arm's charge
    variance.}
  \label{fig:charge}
\end{figure}

\paragraph{Result}
Both the projection and the comparator satisfy the expectation constraint, and the comparator's seed spread on mean charge
is wider than the gap between them, so we claim no ordering there. The difference that
lies outside seed noise is in the last two columns. The
optimum's charge variance exceeds the mask-conditioned data's, \emph{raising} the variance of
the quantity being constrained by $22.1\%$, while the projection lowers it by $81.9\%$. Against the optimum rather than against the unconstrained arm, the
projection's variance is lower by a factor of $6.76$, a gap exceeding three of the comparator's seed
standard deviations. Individual samples follow: the optimum's own mass sits on
constraint-violating sequences $1.42$ times more often than the projection's does.

The mechanism is visible in \cref{fig:charge} and follows from what the tilt is. It is
the KL-closest distribution with the right mean, and reaching a mean of zero from a base
centred at $+2.469$ is cheapest by moving mass symmetrically outward. Nothing in the
objective rewards concentration. Bounding an expected increment at every step, by
contrast, bounds the total in expectation and narrows it as a side effect, because each
step's feasible set is centred on what the remaining budget can still absorb.
%% results/WS1_induced_vs_tilt.json (comparison.excess_kl_of_induced_over_tilt_nats)
The surrogate's $1.025$ nats of excess divergence from the optimum is therefore a trade
rather than a loss, and \cref{si:frontier} reports how that overhead varies with
constraint strength.

\subsection{One checkpoint retargets across the constraint family}
\label{sec:retargeting}

Because the constraint enters \cref{eq:budget} as a target value rather than as a
training condition, a trained checkpoint can be deployed at values it never trained for.
We test this on the crystal task, where the constrained quantity is mass density and the
%% results/WS4_zero_shot_retargeting.json (matrix: 3 checkpoints x 10 targets)
%% results/WS4_retargeting_seed_replication.json (summary.*)
family is ten density targets spanning the data distribution. Three checkpoints, each
trained at one target, are deployed across all ten. Energies quoted here come from a
surrogate (\cref{si:energy}), not density functional theory.

Structural validity is exact at every one of the thirty cells, including targets no
checkpoint trained for and targets beyond the reachable range. The mask is a property of the
grammar, not of the training target, so retargeting cannot break it.

On accuracy, one checkpoint deployed across the sub-ceiling family is level with the
frontier of
%% preregistration/WS4_equivalence_margin.md fixes the margin BEFORE the test;
%% results/WS4_equivalence.json reads it and fails if the file disagrees.
$7$ separately trained models. A difference test is the wrong instrument here. A large
$p$-value is absence of evidence, not evidence of equivalence, and at $n = 7$ paired
targets the smallest attainable two-sided Wilcoxon $p$ is $2/2^7 = 0.0156$ under any
data, which \cref{si:floor} verifies by exhaustive enumeration. We therefore pre-registered an equivalence margin, defined as the retargeted
arm's own mean within-target seed standard deviation, on the reasoning that a difference
below the method's own reproducibility is not actionable and that a margin tied to the
arm's own noise cannot be widened to pass without changing the arm. Two one-sided tests
against that margin of $0.189$ give $p = 0.0023$, and the pre-registered secondary
condition ($p < 0.05$ at a $0.10$ margin) also holds (\cref{si:equivalence}).

The honest form of the claim is parity, not superiority: on a single seed the retargeted
checkpoint appeared \emph{better} than the frontier, but replication across seeds found
it better at only $3$ of $7$ targets. What survives is that one training run is
statistically indistinguishable from seven at one seventh of the training compute, which
is a claim about cost structure rather than about accuracy. \Cref{si:tab:perseed} gives the
per-seed values behind both readings.

\subsection{Where the method stops}
\label{sec:scope}

\paragraph{Steering saturates at a backbone ceiling}
Requested and achieved density track closely at low targets and then flatten
%% figures/fig2_steering_frontier.pdf <- results/FIG2_steering_frontier.json,
%% aggregated seed-wise from results/WS3_steering_frontier_rows.json
(\cref{fig:frontier}). The plateau sits at $10.405$~\unit{\gram\per\centi\metre\cubed}:
requests beyond it are not met, and no amount of additional tilt reaches them, because
the projection can only redistribute mass the backbone already assigns. The constraint
is a reweighting of the model's own proposals, not a source of new modes.

\begin{figure}[t]
  \centering
  \begin{subfigure}[t]{0.48\linewidth}
    %% figures/fig2_steering_frontier.pdf <- results/FIG2_steering_frontier.json
    \kilgraphics{figures/fig2_steering_frontier.pdf}
    \caption{Requested against achieved density, with the ceiling marked.}
    \label{fig:frontier}
  \end{subfigure}\hfill
  \begin{subfigure}[t]{0.48\linewidth}
    %% figures/fig4_infeasibility_vs_target.pdf <- results/WS3_steering_frontier_rows.json
    %% (infeas_init / infeas) and scripts/crystal/analyze_sweep.py's repair factor
    \kilgraphics{figures/fig4_infeasibility_vs_target.pdf}
    \caption{Infeasible steps before and after training, with the repair factor.}
    \label{fig:infeas}
  \end{subfigure}
  \caption{\textbf{The ceiling and its mechanism.} Steering follows the request until the
    backbone runs out of mass to redistribute; the repair the layer performs weakens as the
    request hardens, and where it approaches unity the ceiling binds.}
\end{figure}

\Cref{fig:infeas} gives the mechanism. Before training, a large fraction of decoding
steps present a budget outside the attainable range, and the layer falls back. Training
repairs most of that at easy targets and progressively less at hard ones; where the
repair factor approaches unity, the ceiling binds.

\paragraph{Two constraints compose partially, and asymmetrically}
The second additive constraint is valence electron density, chosen because it is
computable for every element in the split and only partly correlated with mass density
%% results/CP7_second_channel.json (channel_statistics.correlation = 0.725)
($r = 0.725$). The joint solve converges but steers weakly. Across
%% results/WS5_two_constraint_grid.json (summary.*, fixed_valence_columns)
$7$ trained cells, the fraction of the requested valence shift that survives to the
output never exceeds $11.3\%$, and the second constraint's effect is dominated by
budget exhaustion: at most positions the mass budget the first constraint has already
committed leaves nothing for the second to allocate.

The asymmetry is sharper than the aggregate. Asked to \emph{lower} valence density the
constraint moves the right way in $2$ of $2$ cells; asked to \emph{raise} it, in $2$ of
$5$. A running budget is a remaining quantity: spending less of it is always available by
choosing lighter or fewer atoms, while spending more than the prefix has left is not
available once the budget is exhausted. \Cref{si:grid} gives every cell against
its own control. Charge neutrality, the constraint one would most want here, is not
expressible in this class at all: it quantifies existentially over oxidation states,
and collapsing to one charge per element destroys the predicate rather than approximating
it (\cref{si:expressible}).

\paragraph{The constraint class depends on the serialization}
A per-step budget requires that the prefix determine it. Mass density is a ratio, so a
%% results/CP9_serialization_expressibility.json (rows, verdict; 393 structures)
target density fixes a mass budget only once the cell volume is known. Over $393$
structures and $1{,}217$ element positions, a serialization that closes the cell
parameters first makes the budget determinable at $100\%$ of positions, while one that
emits composition first makes it determinable at $0\%$. The constraint class is
therefore a property of the serialization, not of the layer, and a target quantity that
no prefix determines cannot be steered by this mechanism at all.

\subsection{Baselines}
\label{sec:baselines}

\paragraph{Conditional training is more accurate, where it has data}
Training a single model conditioned on the target value is the strongest same-cost
baseline. It is more accurate than the projection wherever the training split has
%% results/WS3_conditional_baseline.json (summary.mae_below_mean, novel_formula_fraction)
support: mean absolute error $0.085$ against $0.196$ below the data mean. It pays for that
in coverage, emitting $5.9\%$ novel formulas against $94.1\%$, and its accuracy degrades
where the conditioning slice thins, precisely where a practitioner asking for an unusual
target is operating.

\paragraph{Best-of-$n$ beats the projection on accuracy, and pays in diversity}
Given enough inference compute, drawing $n$ samples from the unsteered model and keeping
%% results/WS6_best_of_n.json (rows for error, diversity block for unique fraction)
the closest one needs no training and wins on accuracy: it crosses the projection's error
at $n = 16$ (\cref{fig:pareto}). The cost is concentration of the accepted set. Unique
formula fraction falls from $0.824$ at $n=1$ to $0.228$ at $n=64$, while the
projection's $0.782$ does not depend on how hard the target is to hit. Sequential Monte
Carlo, the principled member of the same family, reaches the projection at
%% results/WS6_smc_steering.json (rows; sigma_sweep for the untuned caveat)
$6$ of $10$ targets within $K=64$ particles but is worse than best-of-$n$ at equal
compute by a mean factor of $4.2$; its potential width is untuned, so that figure is a
lower bound on what SMC could achieve rather than a refutation of it.

\begin{figure}[t]
  \centering
  %% figures/fig3_accuracy_diversity_pareto.pdf <- results/FIG3_accuracy_diversity_pareto.json,
  %% built from results/WS6_best_of_n.json and results/WS3_conditional_baseline.json
  \kilgraphics[0.62\linewidth]{figures/fig3_accuracy_diversity_pareto.pdf}
  \caption{\textbf{Accuracy bought by selection is diversity spent.} Mean best-error
    against unique-formula fraction over the sub-ceiling targets. Best-of-$n$ walks down
    and to the left as $n$ grows, trading output variety for accuracy; the projection sits
    off the curve at one forward pass per output.}
  \label{fig:pareto}
\end{figure}

The honest reading is that the projection's advantage is not accuracy per target. It is
that the cost is paid once at training and never again at inference, while the
inference-time family pays per output forever and concentrates its outputs to do so.

%%%%%%%%%%%%%%%%%%%%%%%%%%%%%%%%%%%%%%%%%%%%%%%%%%%%%%%%%%%%%%%%%%%%%%%%%%%%%%%%
